\section{时间演化算符}
\subsection{定义与基本性质}
把初态到任意时刻的线性映射记为
\begin{equation}
\Ket{\psi\left(t\right)}=\hat U\left(t,t_0\right)\Ket{\psi\left(t_0\right)}.
\label{eq:evolution-definition}
\end{equation}
概率守恒要求内积保持不变，故
\begin{equation}
\hat U^\dagger\left(t,t_0\right)\hat U\left(t,t_0\right)=\hat I,
\end{equation}
即 $\hat U\left(t,t_0\right)$ 是酉算符。唯一性与连续演化还给出
\begin{align}
\hat U\left(t_0,t_0\right)&=\hat I,\\
\hat U\left(t_2,t_0\right)&=\hat U\left(t_2,t_1\right)\hat U\left(t_1,t_0\right),\\
\hat U^{-1}\left(t,t_0\right)&=\hat U\left(t_0,t\right)=\hat U^\dagger\left(t,t_0\right).
\end{align}
把式\eqref{eq:evolution-definition}代入 Schr\"odinger 方程，得到演化算符满足
\begin{equation}
\boxed{\mathrm{i}\hbar\frac{\partial}{\partial t}\hat U\left(t,t_0\right)
=\hat H\left(t\right)\hat U\left(t,t_0\right),
\qquad \hat U\left(t_0,t_0\right)=\hat I.}
\label{eq:evolution-equation}
\end{equation}
对第二个时间变量求导则有
\begin{equation}
\mathrm{i}\hbar\frac{\partial}{\partial t_0}\hat U\left(t,t_0\right)
=-\hat U\left(t,t_0\right)\hat H\left(t_0\right).
\end{equation}

\subsection{Dyson 级数与时间排序}
对式\eqref{eq:evolution-equation}从 $t_0$ 到 $t$ 积分，得到积分方程
\begin{equation}
\hat U\left(t,t_0\right)=\hat I-\frac{\mathrm{i}}{\hbar}
\int_{t_0}^{t}\hat H\left(t_1\right)\hat U\left(t_1,t_0\right)\,\mathrm{d}t_1.
\end{equation}
逐次代入右侧的 $\hat U$，得到 Dyson 级数
\begin{align}
\hat U\left(t,t_0\right)
&=\hat I+\sum_{n=1}^{\infty}\left(-\frac{\mathrm{i}}{\hbar}\right)^n
\int_{t_0}^{t}\mathrm{d}t_1\int_{t_0}^{t_1}\mathrm{d}t_2\cdots
\int_{t_0}^{t_{n-1}}\mathrm{d}t_n
\hat H\left(t_1\right)\hat H\left(t_2\right)\cdots\hat H\left(t_n\right).
\label{eq:dyson-series}
\end{align}
积分区域自动满足 $t_1\geq t_2\geq\cdots\geq t_n$，即较晚时刻的算符位于左侧。

定义时间排序算符 $\mathcal T$：对两个算符，
\begin{equation}
\mathcal T\left[\hat H\left(t_1\right)\hat H\left(t_2\right)\right]
=\theta\left(t_1-t_2\right)\hat H\left(t_1\right)\hat H\left(t_2\right)
+\theta\left(t_2-t_1\right)\hat H\left(t_2\right)\hat H\left(t_1\right).
\end{equation}
于是式\eqref{eq:dyson-series}可简写为
\begin{equation}
\boxed{\hat U\left(t,t_0\right)=
\mathcal T\exp\left[-\frac{\mathrm{i}}{\hbar}
\int_{t_0}^{t}\hat H\left(t'\right)\,\mathrm{d}t'\right].}
\label{eq:time-ordered-U}
\end{equation}
时间排序不能一般地省略，因为不同时刻的 Hamilton 算符可能不对易。

\begin{example}
考虑
\begin{equation}
\hat H\left(t\right)=\frac{\hat p^2}{2m}-A\sin\left(\omega t\right)\hat x.
\end{equation}
利用 $\left[\hat x,\hat p^2\right]=2\mathrm{i}\hbar\hat p$，可得
\begin{equation}
\left[\hat H\left(t_1\right),\hat H\left(t_2\right)\right]
=\frac{\mathrm{i}\hbar A}{m}
\left[\sin\left(\omega t_2\right)-\sin\left(\omega t_1\right)\right]\hat p.
\end{equation}
一般情形下该对易子不为零，所以不能把 $\hat U$ 当作普通指数直接计算。
\end{example}

\subsection{可去掉时间排序的情形}
若任意 $t_1,t_2$ 都满足
\begin{equation}
\left[\hat H\left(t_1\right),\hat H\left(t_2\right)\right]=0,
\end{equation}
则各时刻算符的次序无关，从而
\begin{equation}
\hat U\left(t,t_0\right)=
\exp\left[-\frac{\mathrm{i}}{\hbar}\int_{t_0}^{t}\hat H\left(t'\right)\,\mathrm{d}t'\right].
\end{equation}
特别地，若 $\hat H$ 不含时，
\begin{equation}
\boxed{\hat U\left(t,t_0\right)=
\mathrm{e}^{-\mathrm{i}\hat H\left(t-t_0\right)/\hbar}.}
\end{equation}
若 $\hat H=\sum_n E_n\Ket{n}\Bra{n}$，则
\begin{equation}
\hat U\left(t,t_0\right)=
\sum_n\mathrm{e}^{-\mathrm{i}E_n\left(t-t_0\right)/\hbar}\Ket{n}\Bra{n},
\end{equation}
连续谱部分应相应改为积分。
